\section{Discussion}
\label{sec:discussion}

\paragraph{When is cosine enough?} On our evidence, rarely as a decision rule and never as a guarantee. Cosine detects one kind of error with some reliability---content that has no counterpart in the evidence---the most common annotated hallucination type in RAGTruth. It does not detect changed numbers, dates, causal connectives, directions or negations: for these, its paired detection is close to zero, and in more than a third of matched triplets the edit scores at least as high as a paraphrase of the same sentence. In its best configuration it ranks natural answers no better than token overlap with the evidence, and inside a feature model it adds no information beyond lexical features. Its one practical virtue is stability: it rarely rejects a harmless rewording, so a cosine gate that \emph{accepts} only very high-similarity answers wastes little. It is also cheap enough to compute everywhere. The defensible uses are therefore narrow: a first-pass filter for answers that are near-copies of their evidence, a feature among others, or a monitoring signal for gross drift---not a factuality check.

\paragraph{Score movement is not detection.} Every method in our study lowers its score for most single-fact flips, and a study that reported only score drops would conclude that all of them are sensitive. What matters in deployment is whether a fixed threshold changes its decision, and against that criterion detection rates are several times lower (for cosine, \CfCosDrop\% score drops but \CfCosDetect\% paired detection). The inversion rate makes the reason visible: cosine's response to a factual edit is the same size as its response to a paraphrase, so the threshold cannot separate them. We recommend that perturbation studies of factuality metrics report a matched paraphrase control and a decision-level metric alongside score drops.

\paragraph{NLI is the most fact-sensitive cheap method, but not the best ranker.} A zero-shot NLI scorer has the lowest inversion rate by a wide margin and the most uniform detection across error types, yet on natural RAGTruth data it is no better than cosine at its calibrated threshold. Its sensitivity is diluted by sentence averaging in long answers, and its thresholds transfer no better than anyone else's. Trained grounding detectors \citep{tang2024minicheck,zha2023alignscore} are the obvious next step and may change this picture.

\paragraph{Thresholds are the weak point.} The most consistent result of the study is that a calibrated false-acceptance budget is a property of a threshold \emph{and} a distribution. Every evaluator, including the learned risk model, stayed near a 5\% budget in-domain and broke it in several other domains, by up to an order of magnitude; within one benchmark, answer length alone was enough. Any deployment that uses a fixed similarity threshold should recalibrate it per domain on labelled data; the acceptance rate alone is not a reliable alarm, because in several failing transfer cells it moved by only a few points (Table~\ref{tab:transfermatrix}).

\paragraph{Cascades help only as much as the judge.} A cheap risk score resolves more examples than cosine within an error budget, but what that buys depends on the judge. With the 3B judge the cascade mainly \emph{protected} the system from the judge; with the 14B judge it saved a quarter to two fifths of judge calls at an accuracy cost our subset could not pin down. Because the 14B judge was strict on some tasks and lenient on summaries, a judge can fail in either direction, and it needs its own validation and calibration on the target distribution before a cascade is built around it. The random-escalation control is essential for reading these results: only the risk model's margin made targeted escalation clearly better than random.

\section{Ablations and Error Analysis}
\label{sec:ablation}

\begin{table}[t]
\centering\small
\setlength{\tabcolsep}{2.5pt}
\resizebox{\columnwidth}{!}{\begin{tabular}{lrlrrrr}
\toprule
 & \multicolumn{2}{c}{RAGTruth test AUC} & FaithB. & HaluEval & \multicolumn{2}{c}{Counterfactual} \\
Feature set & AUC & $\Delta$ vs.\ full [95\% CI] & AUC & AUC & drop\,\% & AUC \\
\midrule
All groups (full) & .817 &  & .637 & .977 & 79.4 & .606 \\
\midrule
$-$ cosine & .818 & .000 [-.002, .003] & .644 & .983 & 79.4 & .612 \\
$-$ lexical & .794 & -.023 [-.033, -.014] & .583 & .693 & 89.8 & .639 \\
$-$ number & .812 & -.005 [-.009, -.002] & .633 & .976 & 79.3 & .592 \\
$-$ entity & .810 & -.008 [-.012, -.003] & .633 & .974 & 79.0 & .573 \\
$-$ negation & .816 & -.001 [-.002, .000] & .642 & .978 & 93.1 & .608 \\
$-$ length & .808 & -.010 [-.015, -.005] & .644 & .965 & 76.7 & .607 \\
\midrule
only cosine & .746 &  & .535 & .244 & 83.4 & .580 \\
only lexical & .789 &  & .637 & .973 & 88.6 & .551 \\
only number & .565 &  & .504 & .560 & 22.3 & .597 \\
only entity & .663 &  & .581 & .719 & 39.3 & .652 \\
only negation & .547 &  & .580 & .533 & 26.5 & .563 \\
only length & .661 &  & .597 & .975 & 47.8 & .555 \\
\bottomrule
\end{tabular}
}
\caption{Feature-group ablation of the full risk model (logistic regression, RAGTruth training split). $\Delta$: AUC change on the RAGTruth test split when the group is removed, with a 95\% cluster-bootstrap interval. ``Counterfactual'': share of flips that lower the support score, and AUC for separating originals from flips.}
\label{tab:ablation}
\end{table}

\paragraph{Cosine adds nothing on top of lexical grounding.} Removing the three cosine features from the full risk model leaves its RAGTruth test AUC unchanged (\AblMinusCosineLoss\ \AblMinusCosineLossCi; Table~\ref{tab:ablation}), whereas removing the lexical group costs \AblMinusLexicalLoss\ \AblMinusLexicalLossCi\ and is the largest loss. On its own the lexical group (AUC \AblOnlyLexicalAuc) beats the cosine group (\AblOnlyCosineAuc). The number, entity and length groups each contribute a small but reliable amount; the negation group contributes nothing measurable on natural data. A gradient-boosted model on the same features gives similar AUCs (RT-QA \GbmQaAuc\ vs.\ \LrQaAuc\ for logistic regression), so the conclusion does not depend on the linear form.

\paragraph{What the risk model learns.} The largest standardised coefficients (Table~\ref{tab:coef}) are answer length (longer answers are riskier) and token precision; MiniLM document cosine enters with a \emph{positive} risk coefficient once lexical features are present, i.e.\ it is used as a correction term rather than as evidence of support. Length is a legitimate signal on RAGTruth---longer answers have more opportunities to add unsupported material---but it is also what makes the model's thresholds fragile under length shift (Section~\ref{sec:res-transfer}) and what lets it reach near-perfect AUC on HaluEval-QA for reasons unrelated to grounding.

\paragraph{Natural versus counterfactual sensitivity.} Feature groups that matter on natural data and on counterfactuals are different. Dropping the lexical group lowers natural AUC but \emph{raises} the share of flips that lower the score (Table~\ref{tab:ablation}), and the entity group alone separates originals from flips better (counterfactual AUC \AblOnlyEntityCfAuc) than the full model (\AblFullCfAuc). Natural benchmarks reward features that track the prevalent kind of hallucination in their generators---mostly added, unsupported content---and are nearly silent about value substitutions inside an otherwise grounded sentence, which is exactly where cosine fails.

\section{Limitations and Threats to Validity}
\label{sec:limitations}

\paragraph{Judge scale.} Both LLM judges were run on fixed, source-clustered subsets rather than on full test sets, because a CPU-only machine evaluates a 14B model at seconds to tens of seconds per call. Judge intervals are correspondingly wide, comparisons between a judge and a cheap evaluator are made only on the judged subset, and cascade system metrics are reported on those subsets. The judges are two open-weight models with one main rubric; our results say nothing about frontier API judges, which may behave differently. The prompt-sensitivity check covers one alternative rubric for the small judge only.

\paragraph{Counterfactual labels.} The flips are produced by rules and checked by rules: the replaced value must occur in the evidence and the replacement must not. Some flips will nevertheless be wrong (an entity replacement that the evidence happens to support under another name, a negation that yields an awkward but still entailed sentence), and some paraphrases may change meaning despite the filters. A 160-pair audit sheet is released but no human has reviewed it, so every counterfactual number is a measurement against rule-generated labels. Edits are single-sentence and single-fact; they are cleaner than natural hallucinations and should be read as a diagnostic of sensitivity, not as an estimate of field error rates.

\paragraph{Datasets.} RAGTruth responses come from 2023-era generators, FaithBench is small and deliberately hard (it collects summaries on which earlier detectors disagree), HaluEval-QA labels are synthetic and its hallucinated answers differ from correct ones in length and style, and EVOUNA-NQ measures agreement with short reference answers. FactBench was not used. The transfer matrix therefore spans five domains with very different base rates; a FAR exceedance on one of them may reflect label construction as much as evaluator behaviour.

\paragraph{Evaluator coverage.} We use one NLI checkpoint in a zero-shot SummaC-style configuration and do not evaluate trained grounding detectors such as MiniCheck or AlignScore \citep{tang2024minicheck,zha2023alignscore}, which may close much of the gap to the judges at moderate cost. A third embedding model planned in advance was not run (Appendix~\ref{app:deviations}). The risk estimator is a linear model on hand-built features fitted on RAGTruth training data; it inherits RAGTruth's notion of a hallucination.

\paragraph{Statistics.} Cluster bootstrap intervals treat source documents as independent, and the Holm correction covers only the preregistered primary family; the many secondary comparisons in the tables are descriptive. Calibration-time guarantees for the cascade hold only for data exchangeable with the calibration split.

\paragraph{Cost.} Latency was measured on one laptop CPU with batch sizes chosen for that machine. Absolute times, and probably the ratios between methods, will differ on a GPU or with a hosted API; we did not measure either. No commercial API was used, so no dollar costs are reported.


\section{Conclusion}
\label{sec:conclusion}

Embedding cosine is rarely enough to judge factuality. It ranks faithfulness no better than lexical overlap, adds nothing to a feature model, is weaker than reference containment on correctness, and responds to a single changed fact about as much as to a harmless paraphrase, so thresholds fitted on natural data rarely turn an accepted answer into a rejected one. NLI is more sensitive to the fact itself, and a cheap risk score resolves more examples within an error budget, but no threshold survived a change of domain. LLM judges were not a safe default either. A 3B judge failed on long evidence. A 14B judge was the only evaluator that tracked single-fact edits, but it was lenient on summaries and cost \CostJudgeStrongRatio\ times as much as cosine. We recommend treating cosine as one weak feature, testing metrics with matched paraphrase controls and decision-level metrics, and recalibrating thresholds and judges wherever the data change. Code, configurations, raw judge outputs and analysis outputs are available at \CodeURL.
